\documentclass[a4paper,12pt]{article}
\usepackage{utils}

\title{ Questions ouvertes }
\date{\today}

\begin{document}
\maketitle
\tableofcontents
\clearpage

\section{\today}

\subsection{Ce qui à été fait}

Un programme fait en python pour tester en partant d'un chiffre avec un multiple donné \textbf{(Voir annexe 1)}

Le nombre de chiffres dépend du chiffre multipliant de départ.
\begin{itemize}[label=--]
    \item pour 2 on a 18 chiffres
    \item pour 3 on a 28 chiffres
    \item pour 4 on a 6 chiffres
    \item pour 5 on a 42 chiffres sauf 7 x 5
    \item pour 6 on a 58 chiffres
    \item pour 7 on a 22 chiffres
    \item pour 8 on a 13 chiffres
    \item pour 9 on a 44 chiffres
    \item pour 10 on a 2 chiffres
\end{itemize}

\subsection{Les nouvelles questions}

Pourquoi le nombre de chiffres dépend du chiffre multipliant de départ ?

\clearpage

\section{\today}

\subsection{Ce qui à été fait}

amélioration du programme, qui prend maintenant en compte les multipliant supérieurs à 9 et peut afficher plusieurs nombres possibles avec gradient sur les chiffres.

\begin{tcolorbox}[
        sharp corners,
        colback=white,
        colframe=black,
        leftrule=1mm,
        toprule=-1mm, bottomrule=-1mm, rightrule=-1mm,
        top=0mm, left=1mm
    ]
    Soient 
    \begin{enumerate}[label=--]
        \item $x$ le nombre tournant recherché,
        \item $d$ le chiffre des unités de $x$,
        \item $n$ le nombre total de chiffres de $x$,
        \item $m$ le multipliant.
    \end{enumerate}

    On pose $y := \dfrac{x - d}{10} \implies x = 10y + d$ le nombre obtenu en retirant le chiffre des unités de $x$ et $x' = d \times 10^{n - 1} + y$ le nombre obtenu en mettant le chiffre des unités de $x$ au début de $x$.

    D'après notre sujet de recherche, on doit trouver $x$ ou $x'$ tel que $x' = m \times x$, donc:
    \begin{align*}
        x' = m \times x &\implies d \times 10^{n - 1} + y = m \times (10y + d) \\
        &\implies y - 10my = md - d \times 10^{n - 1} \\
        &\implies y(1 - 10m) = d(m - 10^{n - 1}) \\
        &\implies y = \frac{d(m - 10^{n - 1})}{1 - 10m} \\
    \end{align*} 

    Sauf qu'on ne connait pas $n$, on essaye de chercher des valeurs de $n$ telles que $y$ soit un entier positif. On en déduit ensuite $x = 10y + d$.
\end{tcolorbox}


\subsection{Les nouvelles questions}

\clearpage

\section{\today}

\subsection{Ce qui à été fait}

\begin{definition}
    Soit $n, m \in \N$ avec $m \neq 0$.
    Un nombre $n$ est dit $m$-parasite si \text{$n \times m = d10^{L - 1} + \dfrac{n - d}{10}$} où $d$ est le chiffre des unités de $n$ et $L$ sa taille (nombre de chiffres).
\end{definition}

\begin{exemple}
    $$
    128205 \times 4 = 512820
    $$
\end{exemple}

Voir annexe 3 la démonstration

\subsection{Les nouvelles questions}

\clearpage

\section{\today}

\subsection{Ce qui à été fait}

\begin{itemize}
    \item Changement de base.
    \item Est-ce que pour toute longueurs il existe un nombre tournant de cette longueur.
    \item On a trouvé un nombre qui a des rotations différentes suivant la multiplication : $142 857$
    \begin{itemize}
        \item $142 857 \times 2 = 285 714$
        \item $142 857 \times 3 = 428 571$
        \item $142 857 \times 4 = 571 428$
        \item $142 857 \times 5 = 714 285$
        \item $142 857 \times 6 = 857 142$
    \end{itemize}
\end{itemize}

\subsection{Les nouvelles questions}

\clearpage

\section{\today}

\subsection{Ce qui à été fait}

\begin{itemize}
    \item rotation quelconque
    \item début de préparation oral
    \item réponse à la question si toute longueur il existe un nombre tournant
    \item À faire chercher les bases les plus simples
\end{itemize}

\subsection{Les nouvelles questions}

\clearpage

\section{Annexes}

\subsection*{1}

\begin{lstlisting}
def until_turned(first_value, mul, limit=100):
    value = [first_value]
    result = []
    retenu = 0
    for i in range(limit):
        if result and retenu == 0 and result[-1] == first_value:
            break
        num = (value[-1] * mul + retenu) % 10
        retenu_tmp = (value[-1] * mul + retenu) // 10
        result.append(num)
        value.append(num)
        retenu = retenu_tmp
    return (''.join(map(str, value[::-1][1:])), len(value) - 1, ''.join(map(str, result[::-1])), len(result))

mul = 6
for i in range(1, 10):
    print(f"pour {i} avec mul = {mul} on a {until_turned(i, mul)}")
\end{lstlisting}

\subsection*{3}

\begin{lemme}{}{existence}
    Soit $n, p \in \N$

    Supposons que $p$ est inversible dans $\Z/n\Z$.

    Alors $ord(p)$ existe dans $\Z/n\Z$
\end{lemme}

\begin{demonstration}
    Comme $\Z/n\Z$ est fini, on a forcément $\exists k, l \in \N, k > l$ tel que
    \begin{align*}        
        p^k \equiv p^l [n] \implies p^k (p^{-1})^l \equiv p^l (p^{-1})^l [n] \implies p^{k-l} \equiv 1 [n]
    \end{align*}
\end{demonstration}

\begin{theoreme}{Caractérisation de la taille d'un nombre $m$-parasite}{}
    Soit $m \in \N^*$, $x$ un nombre $m$-parasite de chiffre des unités $d$ et $L$ sa taille (nombre de chiffres).

    Alors $L$ peut prendre des valeurs dans $\{ord(10) \times k \mid k \in \N^*\}$ dans $\Z/(\dfrac{10m - 1}{pgcd(d, 10m - 1)})\Z$.
\end{theoreme}

\begin{demonstration}
    On pose $y := \dfrac{x - d}{10}$ le nombre obtenu en retirant le chiffre des unités $d$ de $x$.

    Ainsi, on peut définir $x = 10y + d$.

    par définition d'un nombre $m$-parasite, on a
    \begin{align*}
        m \times x = d10^{L - 1} + \dfrac{x - d}{10} &\implies m(10y + d) = d10^{L - 1} + y \\
        &\implies dm - d 10^{L-1} = y - 10my \\
        &\implies y(1 - 10m) = d (m - 10^{L-1}) \\
        &\implies y = d \dfrac{m - 10^{L - 1}}{1 - 10m} = d \dfrac{10^{L - 1} - m}{10m - 1}
    \end{align*}

    Comme $y$ est un entier, le numérateur doit être divisible par le dénominateur.

    Donc
    \begin{align*}
        \exists k \in \N, &d(10^{L - 1} - m) = k(10m - 1) \quad (*) \\
        \implies& d(10^{L - 1} - m) \equiv 0 [10m - 1] \\
        \implies& d 10^{L - 1} \equiv dm [10m - 1]
    \end{align*}

    $d$ est inversible dans $\Z/(10m - 1)\Z$ si et seulement si \text{$pgcd(d, 10m - 1) = 1$}.

    Ainsi on a deux cas :
    \begin{description}
        \item[si $pgcd(d, 10m - 1) = 1$ :] on peut simplifier par $d$ et on obtient
        \begin{align*}
            &10^{L - 1} \equiv m [10m - 1] \\
            \implies& 10^L \equiv 10m [10m - 1] \\
            \implies& 10^L \equiv 1 [10m - 1]
        \end{align*}

        De plus $m$ est l'inverse de $10$ dans $\Z/(10m - 1)\Z$.

        D'apres le lemme \ref{lemme:existence}, $ord(10)$ ainsi que ses multiples existe bien dans $\Z/(10m - 1)\Z$ et résoud notre équation.

        Donc $L \in \{ord(10) \times k \mid k \in \N^*\}$ dans $\Z/(10m - 1)\Z$.
    
        \item[si $pgcd(d, 10m - 1) \neq 1$ :] on peut poser $gcd = pgcd(d, 10m - 1)$,
        ainsi avec $(*)$ on a
        
        \begin{align*}
            \exists k \in \N, &\dfrac{d(10^{L - 1} - m)}{gcd} = k\dfrac{(10m - 1)}{gcd} \\
            \implies& \dfrac{d}{gcd} 10^{L - 1} \equiv \dfrac{d}{gcd} m [\dfrac{10m - 1}{gcd}]
        \end{align*}

        Cette fois $pgcd(\dfrac{d}{gcd}, \dfrac{10m - 1}{gcd}) = 1$ donc $\dfrac{d}{gcd}$ est inversible, d'où
        \begin{align*}
            &10^{L - 1} \equiv m [\dfrac{10m - 1}{gcd}] \\
            \implies& 10^L \equiv 10m [\dfrac{10m - 1}{gcd}] \\
            \implies& 10^L \equiv 1 [\dfrac{10m - 1}{gcd}]
        \end{align*}

        car
        \begin{align*}
            10m - 1 = gcd \dfrac{10m - 1}{gcd} &\implies 10m - 1 \equiv 0 [\dfrac{10m - 1}{gcd}] \\
            &\implies 10m \equiv 1 [\dfrac{10m - 1}{gcd}]
        \end{align*}

        De plus $m$ est l'inverse de $10$ dans $\Z/(\dfrac{10m - 1}{gcd})\Z$

        D'apres le lemme \ref{lemme:existence}, $ord(10)$ ainsi que ses multiples existe bien dans $\Z/(\dfrac{10m - 1}{gcd})\Z$ et résoud notre équation.

        Donc $L \in \{ord(10) \times k \mid k \in \N^*\}$ dans $\Z/(\dfrac{10m - 1}{gcd})\Z$.
    \end{description}

    Comme dans le premier cas $pgcd(d, 10m - 1) = 1$ on peut rassembler les deux cas en un seul

    $L \in \{ord(10) \times k \mid k \in \N^*\}$ dans $\Z/\dfrac{10m - 1}{pgcd(d, 10m - 1)}\Z$.
\end{demonstration}

\end{document}